\documentclass[10pt,a4paper]{amsart}
\usepackage[margin=19mm]{geometry}
\usepackage{amsmath,amssymb,mathtools}
\usepackage[scr=rsfs,cal=euler]{mathalfa}
\usepackage{xfrac}
\usepackage[unicode,hidelinks,bookmarks=false]{hyperref}
\newcommand{\ie}{i.e.\ }
\newcommand{\cf}{cf.\ }
\newcommand{\resp}{respectively\ }
\newcommand{\Subsection}{\subsection}
\providecommand{\nobreakdash}{\nobreak\hbox{-}}
\setlength{\emergencystretch}{2em}
\multlinegap0pt
\usepackage{amssymb}
\usepackage{mathtools}
\usepackage{booktabs}
\usepackage[shortlabels]{enumitem}
\usepackage{tikz}
\usepackage{float}
\newtheorem{lemma}{Lemma}
\title{Stability phenomena in Deligne--Mumford compactifications via Morse theory}
\author{Changjie Chen}
\date{Source: arXiv:2511.13695v2 (CC BY 4.0)}
\begin{document}
\maketitle

\noindent\emph{Source and licence.} Stability phenomena in Deligne--Mumford compactifications via Morse theory. Version arXiv:2511.13695v2 (CC BY 4.0), distributed by arXiv under the Creative Commons Attribution 4.0 International licence, http://creativecommons.org/licenses/by/4.0/, as recorded in the arXiv licence metadata for this exact version and shown on the abstract page https://arxiv.org/abs/2511.13695v2. Full licence notice: \texttt{licenses/arxiv-2511.13695-CC-BY-4.0.txt}.\par\smallskip

\noindent\textit{Editorial scope.} Counted unit: the article's lemma lemma (source0.tex lines 342--344). The article's complete original proof of it is delivered with it (source0.tex lines 346--366, one \texttt{proof} environment of 21 lines). No mathematics is written, repaired or re-derived: the statement and the proof are the article's own lines, in the article's own notation, and no retained line is substituted or re-flowed.  No further source-local context is needed: every cross-reference of every retained line resolves inside the two delivered spans.  Nothing was omitted from any retained span, so the package carries no omission record.  No retained line contains a citation, so the wrapper carries no bibliography and no bibliography entry.  No retained line contains a figure environment, an image-inclusion command or drawing code, so no image is delivered and the package declares no asset dependency.  Editorial formatting only, in addition: an AMS \texttt{amsart} preamble that restates the article's own declarations and macros used by the retained lines, namely the environments \texttt{lemma} on separate counters, together with the standard packages those lines need.  The retained environments print in this document as Lemma 1 in order of appearance, whereas the article prints the same items with the numbering of its own full document.  The complete native source of the article as distributed in the arXiv e-print for this exact version is bundled unchanged as \texttt{sources/189308/source0.tex}.
\begin{lemma}
    Let $f\colon M\to\mathbb R$ be a Morse function on a manifold $M$ equipped with a Riemannian metric, and $N\subset M$ a Morse submanifold of $M$. Then if the closure of a flow line of $f$ intersects $N$, then it is contained in $N$ or hits a critical point in $N$ transversely.
\end{lemma}

\begin{proof}
    Let $\varphi$ be the closure of an $f$-flow line. If $\varphi$ contains a point in $N$ in its interior, then $\varphi$ is contained in $N$ since $N$ is a Morse submanifold of $M$. Suppose $\varphi$ hits a point $p\in N$ from outside $N$, then $p$ lies at the end of $\varphi$ and is a critical point of $f$. Take the orthogonal decomposition of the tangent space as
    \[
    T_pM=T_pN\oplus \nu_pN.
    \]
    Let $(x,y)$ be a compatible local chart around $p$ with $(x,y)(p)=(0,0)$ such that
    \[
    f(x,y)=f(p)+x^TAx+y^TBy+x^TCy+g(x,y)
    \]
    where $A$ and $B$ are nondegenerate matrices, and $\nabla^2 g=0$. Note that
    \[
    \nabla f=(2Ax+Cy,2By+Cx),
    \]
    and therefore, $C=0$ since locally $N=\{y=0\}$.

    Pick $p_0=\varphi(t_0)=(x_0,y_0)$ where $y_0\neq0$, then the flow through $p_0$ is given by
    \[
    \varphi(t-t_0)=(e^{2At}x_0,e^{2Bt}y_0),
    \]
    which is transverse to $N$ if it converges to $p$ as $t\to-\infty$.
\end{proof}

\end{document}
